% Part: applied-modal-logics
% Chapter: epistemic-logic
% Section: introduction

\documentclass[../../../include/open-logic-section]{subfiles}

\begin{document}

\olfileid{aml}{el}{int}

\olsection{పరిచయం}

మోడల్ తర్కం \emph{మోడల్ ప్రతిపాదనలు}, వాటి మధ్యనున్న అనుగమన సంబంధాలను అధ్యయనం చేసినట్లే, జ్ఞానసంబంధ తర్కం \emph{జ్ఞానసంబంధ ప్రతిపాదనలు}, వాటి మధ్యనున్న అనుగమన సంబంధాలను అధ్యయనం చేస్తుంది. మోడల్ కారకాలను సాధ్యత, ఆవశ్యకతల సూచకాలుగా అర్థం చేసుకునే బదులు, జ్ఞానాన్ని, విశ్వాసాన్ని నమూనీకరించడానికి ఏకస్థానిక సంయోజకాలను జ్ఞానసంబంధ లేదా విశ్వాససంబంధ పద్ధతుల్లో అర్థం చేసుకుంటాం. ఉదాహరణకు, కింది వాదనలను వ్యక్తీకరించాలనుకోవచ్చు:

\begin{enumerate}
\item కాల్గరీ ఆల్బర్టాలో ఉందని రిచర్డ్‌కు తెలుసు.
\item సోఫాపై ఒక కుక్క ఉండే అవకాశం ఉందని ఆడ్రీ భావిస్తుంది.
\item ఆడ్రీకి తన తరగతి మంగళవారాల్లో ఉంటుందని తెలుసని రిచర్డ్‌కు తెలుసు.
\item ఒక సంవత్సరంలో 12 నెలలు ఉంటాయని అందరికీ తెలుసు.
\end{enumerate}
సమకాలీన జ్ఞానసంబంధ తర్కానికి మూలంగా తరచూ జాకో హింటిక్కా 1962లో రాసిన \emph{Knowledge and Belief}ను పేర్కొంటారు. సాధ్య ప్రపంచాల అర్థవిజ్ఞానం తర్కశాస్త్రంలో క్రమంగా విస్తృతంగా వాడుకలోకి వస్తున్న కాలంలో ఈ గ్రంథం వెలువడింది. నిజానికి, జ్ఞానసంబంధ తర్కాలు ఇతర మోడల్ తర్కాల అర్థవిజ్ఞాన సాధనాల్లో చాలా వాటిని వాడతాయి; అయితే వాటికి భిన్నమైన అర్థాన్ని ఇస్తాయి. ముఖ్యమైన మార్పు \emph{ప్రాప్యత సంబంధం} దేనిని సూచిస్తుందని మనం తీసుకుంటామనే విషయంలో ఉంటుంది. జ్ఞానసంబంధ తర్కాల్లో అది ఏదో రూపంలోని \emph{జ్ఞానపర సాధ్యత}ను సూచిస్తుంది. మనం నమూనీకరిస్తున్న జ్ఞానసంబంధ భావనను బట్టి ప్రాప్యత సంబంధంపై విధించదలచే షరతులు మారుతాయని చూస్తాం. అనురూపతా సిద్ధాంతానికి జ్ఞానసంబంధ అర్థం ఇచ్చినప్పుడు ఏమి జరుగుతుందో కూడా చూస్తాం. పై ఉదాహరణల్లో రిచర్డ్, ఆడ్రీ అనే ఇద్దరు కర్తలు, ఒక్కొక్కరికి తెలిసిన విషయాల మధ్య సంబంధం ఉన్నాయని గమనించవచ్చు. మనం పరిశీలించే జ్ఞానసంబంధ తర్కాలు ఇటువంటి విషయాలను వ్యక్తీకరించగల బహు-కర్త తర్కాలు. దీనికి భిన్నంగా, ఏక-కర్త జ్ఞానసంబంధ తర్కం ఒక్క వ్యక్తికి తెలిసిన లేదా అతను విశ్వసించే విషయాల గురించి మాత్రమే మాట్లాడుతుంది.

\end{document}
